\section{Prototype}
\label{sec:pat-prototype}

\problemstatement{Create new objects by cloning an existing prototype
rather than constructing from scratch, useful when object creation is
expensive or when the exact concrete type to copy is known only at run
time.}

\begin{patternbox}[When You Reach for It]
The trigger is \emph{``make a copy of an existing configured object,''} or
\emph{``creating one from scratch is costly.''} When you have a
fully-configured instance and want more just like it -- without knowing or
naming its concrete class -- a polymorphic \texttt{clone()} is the answer.
\end{patternbox}

\subsection*{Understanding the pattern}

Each class implements a virtual \texttt{clone()} that returns a copy of
itself as the base type. Client code holds a base pointer and calls
\texttt{clone()}; it gets back a new object of the correct concrete type
without a \texttt{switch} on types or a call to a specific constructor.
This is invaluable when the set of concrete classes is registered at run
time (a palette of prototype shapes the user can duplicate).

\begin{ideabox}[Let the Object Copy Itself]
Because only the object knows its own concrete type, delegate copying to a
virtual \texttt{clone()}. This sidesteps ``which constructor do I call?''
and preserves all the run-time state of the prototype, including fields a
factory would not know to set.
\end{ideabox}

\subsection*{C++ implementation}

\begin{lstlisting}
#include <bits/stdc++.h>
using namespace std;

struct Shape {
    int x = 0, y = 0;
    virtual unique_ptr<Shape> clone() const = 0;   // polymorphic copy
    virtual void draw() const = 0;
    virtual ~Shape() = default;
};

struct Circle : Shape {
    int r = 1;
    unique_ptr<Shape> clone() const override {
        return make_unique<Circle>(*this);          // copy all state
    }
    void draw() const override {
        cout << "Circle r=" << r << " at (" << x << "," << y << ")\n";
    }
};

int main() {
    Circle base; base.r = 5; base.x = 2; base.y = 3;

    unique_ptr<Shape> proto = make_unique<Circle>(base);
    auto copy1 = proto->clone();                    // no concrete type named
    auto copy2 = proto->clone();
    copy1->draw();
    copy2->draw();
    return 0;
}
\end{lstlisting}

\begin{complexitybox}[Trade-offs]
\textbf{Pros.} Copies run-time state exactly; avoids costly
re-construction; adds objects without naming concrete classes.
\textbf{Cons.} Deep vs.\ shallow copy must be handled carefully
(pointers, cycles); every class must implement \texttt{clone()}
correctly. Watch out for objects owning raw pointers.
\end{complexitybox}

\begin{takeawaybox}
Prototype creates by copying a live object through a virtual
\texttt{clone()}, sidestepping constructors and type switches. Reach for
it when you already have a configured instance and want duplicates, or
when construction is expensive. The one hazard is getting deep-copy
semantics right.
\end{takeawaybox}
